\begin{abstract}
Model selection for verifiable tasks should ask how much resource is required
to obtain a verified solution, not only which model has the highest one-shot
accuracy.  We formalize this objective as a reliability-targeted quantile of
verified first-passage resource, measured in a predeclared unit such as time,
tokens, energy, or money.  Simultaneous finite-sample confidence bands provide
a conservative rule for admitting and comparing a finite menu of agent
systems, including an explicit no-winner outcome.  To explain measured
differences, we introduce a batched tool-allocation process and isolate
a restricted reference-combination regime.  In that regime, expected log
conditional first-passage resource separates exactly into relative mean retry
cost, focused competence, signal information about appropriate tool use, and
batched tool-policy mismatch.  Controlled Qwen2.5-Coder
experiments recover the one-dispatch special case with near-unit inverse-share
scaling and decision-accurate accounting while rejecting exact closure; public
benchmark traces further show that reliability targets are often unattained.
Thus verified proper time selects the deployment, whereas the four-term
identity audits why that selection succeeds or fails.
\end{abstract}

\section{Introduction}
\label{sec:introduction}

Which model should be deployed on a new task whose solution can be checked?
Current practice often evaluates candidate systems end to end and deploys the
one with the best aggregate score.  We call this \emph{benchmark racing}.  It
has produced realistic evaluations, from repository repair against executable
tests to machine-learning engineering against held-out criteria
\citep{jimenez2024swebench,chan2024mlebench}.  Yet aggregate accuracy is
insufficient when models differ in serving cost, can be retried, or have
complementary strengths.  A model that is more accurate per attempt need not
reach the first verified solution with the least time, computation, energy,
tokens, or money.  The same aggregate gain can also arise from cheaper attempts,
greater competence, better pre-solution information, or better allocation.
These mechanisms imply different engineering interventions.

Our conceptual foundation is the proper-time framework of Achille and Soatto
\citep{achille2025agents}.  They study verifiable transductive tasks and argue
that eventual correctness under unbounded search does not characterize useful
intelligence.  Indeed, IID retries with success probability \(p>0\) satisfy
\(\Prob(\text{success by }n)=1-(1-p)^n\to1\).  We therefore require a verified
solution within a predeclared resource horizon and reliability target.  Their
path proper time motivates the principle that information is useful insofar as
it buys speed; our deployment quantity is different: a population quantile of
observed verifier-backed first passage.

Our primary contribution is a finite-sample selection rule for that quantile.
It treats distinct task instances as independent observations, permits paired
agent-system runs on the same instance, and either certifies separation or
reports no winner.  A batched tool-allocation model
then connects tool selection to first-passage resource.  Its restricted
reference-combination regime yields the diagnostic identity against an
uninformative, reference-prior-matched baseline,
\begin{equation}
\begin{aligned}
\Delta_{\mathrm{tool}}
={}&\underbrace{\E_{X,Z_0}\log\overline k_{\sigma_0}(X,Z_0)
-\E_{X,Z}\log\overline k_\sigma(X,Z)}_{\text{relative mean trial cost}}
+\underbrace{\E_X\!\left[\log
  \frac{t_0(\sigma_0,X)}{t_0(\sigma,X)}\right]}_{\text{focused competence}}\\[-0.15em]
&+\underbrace{G_Z}_{\text{tool information}}
-\underbrace{\varepsilon_\sigma}_{\text{tool-policy mismatch}}.
\label{eq:intro-four-term}
\end{aligned}
\end{equation}
This identity concerns expected log \emph{conditional mean} resource inside
this regime.  It does not decompose the proper-time quantile and is not used to select
a model.  Direct quantile measurement selects; Equation~\eqref{eq:intro-four-term}
explains the mechanism when its assumptions survive held-out diagnostics.

We make three contributions.  First, we define reliability-targeted verified
proper time and give simultaneous finite-sample CDF, quantile, admissibility,
and winner guarantees for a finite agent-system menu.
Second, we derive the exact mean first-passage resource of a batched tool
policy under complete IID retries.  A transparent reference-combination
specialization yields a mismatch that is a sum of per-decision tool-policy
divergences.  Third, we test the one-dispatch restriction with
real language models.  In 175,104 held-out executions over 768 controlled
tasks, the inverse-share exponent is \(0.987\), the four-term residual RMS is
\(0.015\) nat, and model choice agrees with the held-out oracle in all six conditions.
A strict omnibus test rejects exact closure, so the result is
decision-accurate approximate accounting rather than a universal equality.

\section{Related Work}
\label{sec:related}

Classical algorithm selection maps instance features to a solver; empirical
performance models and SATzilla show that runtime-optimal solvers can be
instance dependent~\citep{rice1976algorithm,hutter2014runtime,xu2012satzilla,
kotthoff2016algorithm}.  Racing and best-arm procedures reduce evaluation cost
when identifying a winner~\citep{maron1993hoeffding,maron1997racing,
kaufmann2016bestarm}.  Run2Survive models censored runtime distributions and
risk-sensitive selection~\citep{tornede2020run2survive}, while preferences over
runtime distributions formalize objectives beyond the mean
\citep{graham2023runtime}.  Our target is likewise distributional, but
specialized to verifier-backed first passage and simultaneous quantile
selection.

LLM cascades and routers learn cost--quality decisions over heterogeneous model
pools~\citep{chen2023frugalgpt,ong2024routellm,hu2024routerbench,
dekoninck2025routing}; test-time scaling further shows that sampling effort and
model size interact~\citep{brown2024monkeys,snell2025testtime,wu2025inference,
osca2024}.  We instead freeze the complete agent system and compare its
first-passage distribution.  Finally, the information term in our restricted
identity is the standard reduction in log loss supplied by side information
\citep{jiao2014logloss}.  Our contribution is not a new characterization of
mutual information, but its placement beside measured cost and competence in a
falsifiable first-passage accounting.

\section{Verifier-Backed Agent-System Selection}
\label{sec:problem}

\begin{definition}[Verifiable task]
\label{def:task-family}
A \emph{verifiable task} is a tuple
\(\mathrm{T}=(\mathrm{X},\mathrm{Y},V,\mu)\).  Each \(x\in\mathrm{X}\) is a
concrete instance for which an artifact \(y\in\mathrm{Y}\) must be constructed.
The distribution \(\mu\) describes future instances, and the predeclared
external verifier \(V:\mathrm{X}\times\mathrm{Y}\to\{0,1\}\) supplies the
common success event.  A \emph{resource measure} assigns every partial
execution a nonnegative additive coordinate \(B\), measured in isolated
GPU-seconds, joules, tokens, dollars, or a predeclared scalarization.
\end{definition}

In this work, an \emph{agent system} \(\sigma\in\Sigma\) consists of a backbone model \(M_\sigma\), its prompting and decoding strategies, available tools,
memory and state, tool and sub-agent orchestration.  Given \(x\), it produces a generally random execution and artifact \(y(B;x)\) using a resource \(B\), which could be tokens, time or money.  The verifier \(V\) is external to \(\sigma\).

\begin{definition}[Verified first passage]
\label{def:first-passage}
The verified first-passage resource is
\begin{equation}
B_\sigma^{\mathrm{hit}}(x):=
\begin{cases}
\inf\{B\ge0:V(x,y(B;x))=1\},&\text{if verification is reached},\\
\infty,&\text{otherwise}.
\end{cases}
\label{eq:first-passage}
\end{equation}
\end{definition}

\begin{definition}[Reliability-targeted verified proper time]
\label{def:proper-time}
For failure tolerance \(\delta\in(0,1)\), define
\begin{equation}
T_\sigma(\delta):=\inf\{B\ge0:\Prob[B_\sigma^{\mathrm{hit}}(x)\le B]\ge1-\delta\}.
\label{eq:proper-time-system}
\end{equation}
This population \((1-\delta)\)-quantile is the smallest resource by which a
verified artifact is produced with the target reliability.
The term \emph{statistically certified} is reserved for the finite-sample bounds
in Section~\ref{sec:finite-sample}.  For finite positive quantiles, define the
proper-time log-speedup
\(\Delta_T(\sigma_0,\sigma;\delta)=
\log T_{\sigma_0}(\delta)-\log T_\sigma(\delta)\).
\end{definition}

\begin{problem}[Reliability-first agent-system selection]
\label{prob:model-selection}
Given a finite predeclared menu \(\Sigma\), common verifier and resource
unit, tolerance \(\delta\), and horizon \(B\), certify which systems satisfy
\(\Prob(B_\sigma^{\mathrm{hit}}\le B)\ge1-\delta\), and select the admissible
system with smallest \(T_\sigma(\delta)\).
\end{problem}

\section{Finite-Sample Certification}
\label{sec:finite-sample}

The proper time in Definition~\ref{def:proper-time} is a population quantity,
whereas an experiment observes finitely many task instances.  For each selected agent system \(\sigma\), run the predeclared first-passage evaluation once on every instance, and record
\(B_i^\sigma\in[0,\infty]\), with \(B_i^\sigma=\infty\)
when no verified artifact is produced by the horizon.  Define
\begin{equation}
\widehat F_\sigma(b):=\frac{1}{n}\sum_{i=1}^{n}
\mathbf 1\{B_i^\sigma\le b\},
\qquad
F_\sigma(b):=\Prob(B_\sigma^{\mathrm{hit}}\le b).
\label{eq:empirical-cdf}
\end{equation}
The independent units are the task instances.  Runs of different systems on
the same instance may be paired and dependent.

\begin{theorem}[Instance-level DKW band]
\label{thm:dkw-system}
For a fixed agent system \(\sigma\), the Dvoretzky--Kiefer--Wolfowitz inequality
with Massart's sharp constant~\citep{massart1990tight} gives
\begin{equation}
\Prob\!\left(
\sup_{b\ge0}|\widehat F_\sigma(b)-F_\sigma(b)|
>\sqrt{\frac{\log(2/\alpha)}{2n}}
\right)\le\alpha.
\label{eq:dkw-system}
\end{equation}
\end{theorem}

\begin{corollary}[Simultaneous proper-time certificate]
\label{thm:finite-sample}
For a finite candidate menu \(\Sigma\), let
\begin{equation}
\begin{aligned}
\eta_n(\alpha)&:=\sqrt{\frac{\log(2|\Sigma|/\alpha)}{2n}},\\
L_\sigma(b)&:=\max\{0,\widehat F_\sigma(b)-\eta_n(\alpha)\},\qquad
U_\sigma(b):=\min\{1,\widehat F_\sigma(b)+\eta_n(\alpha)\}.
\end{aligned}
\label{eq:simultaneous-radius}
\end{equation}
With probability at least \(1-\alpha\),
\(L_\sigma(b)\le F_\sigma(b)\le U_\sigma(b)\) simultaneously for every
\(\sigma\in\Sigma\) and every \(b\ge0\).  For target reliability
\(1-\delta\), define
\begin{equation}
T_\sigma^-:=\inf\{b\ge0:U_\sigma(b)\ge1-\delta\},
\qquad
T_\sigma^+:=\inf\{b\ge0:L_\sigma(b)\ge1-\delta\},
\label{eq:quantile-band}
\end{equation}
with \(\inf\varnothing=\infty\).  On the same event,
\begin{equation}
T_\sigma^-\le T_\sigma(\delta)\le T_\sigma^+.
\label{eq:proper-time-band}
\end{equation}
At horizon \(B\), system \(\sigma\) is certified to meet the target when
\(L_\sigma(B)\ge1-\delta\).  A certified system \(\widehat\sigma\) is the
unique proper-time winner when
\begin{equation}
T_{\widehat\sigma}^{+}<
\min_{\sigma\ne\widehat\sigma}T_\sigma^{-};
\label{eq:certified-winner}
\end{equation}
otherwise the certificate returns \emph{no winner}.
\end{corollary}

For inference about one fixed system, set \(|\Sigma|=1\).  The simultaneous
correction is needed only when the same experiment supports a comparison or a
winner claim across multiple systems.

\section{Tool Allocation and Four-Term Accounting}
\label{sec:tool-accounting}

\begin{definition}[To be decided]
\label{def:reference-tools}
Fix a finite available tool set \(\mathcal A\): the external callable
capabilities available to the agentic LLM inside agent system \(\sigma\).
Elements of \(\mathcal A\) may be functions, APIs, retrievers, code executors,
browsers, or orchestration primitives such as spawning a sub-agent. Assume that
every task instance \(x\) admits at least one finite combination of tool calls
that can produce a verified artifact. We represent such a combination as a
directed acyclic \emph{tool-call graph}: nodes are tool calls, and directed
edges encode data or control dependencies between calls. A call with incoming
edges may use the outputs of its parent calls, while calls with no dependency
relation may be executed in parallel. The reference graph is not assumed to be
unique; other finite tool-call graphs may also lead to a verified solution.


\end{definition}

\begin{definition}[Tool-allocation policy]
\label{def:tool-policy}
Fix an agent system with available tool set \(\mathcal A\). During one complete
attempt, the agent produces a tool-use trace
\(\tau=(C_1,\ldots,C_L)\), where each \(C_k\subseteq\mathcal A\) is the batch of
tool calls dispatched at step \(k\). The empty batch \(C_k=\varnothing\) means
that no external tool is called at that step; batches with multiple elements
represent parallel tool use.

Given a pre-solution signal \(z\) and the previous tool-call prefix
\(\tau_{<k}=(C_1,\ldots,C_{k-1})\), the next batch is drawn from
\[
C_k\sim q(\cdot\mid z,\tau_{<k}),
\qquad
C_k\in 2^{\mathcal A}.
\]
Thus \(q\) is the agent's tool-allocation policy: it specifies how the system
allocates execution across available tools as the attempt unfolds. If
\(\tau^\star(x)=(C_1^\star,\ldots,C_L^\star)\) is a reference successful trace
for instance \(x\), then the probability that the agent reproduces that trace is
\[
q^\star(x,z)
=
\prod_{k=1}^{L}
q\!\left(C_k^\star\mid z,\tau^\star_{<k}\right).
\]
\end{definition}

With retries, the verified first-passage budget is the cumulative resource
consumed until the first verified artifact.  For an agent system \(\sigma\), task
instance \(x\), and optional pre-solution signal \(Z=z\), if retry \(r\) produces
artifact \(Y_r\) at cost \(k_r>0\), then
\[
B_\sigma^{\mathrm{hit}}(x,z)
=
c_\sigma(z)+\sum_{r=1}^{R_\sigma(x,z)} k_r,
\qquad
R_\sigma(x,z):=\min\{r\ge1:V(x,Y_r)=1\}.
\]
Here \(c_\sigma(z)\ge0\) is the incremental cost of obtaining the signal; set
\(c_\sigma(z)=0\) when the signal is supplied with the task or its cost is
already included in the retry costs.  Thus an agent system improves proper time
by reducing the cost of its attempts, increasing the probability that attempts
verify, or using pre-solution information to do both.

\subsection{Reference-combination regime}

Fix an agent
system, and let \(\tau^\star(x)\) be one successful tool-use trace for instance
\(x\), obtained either from benchmark-provided tool traces or from a teacher
system that has solved the instance at least once and therefore provides an
oracle reference trace. Verifiable tasks, such as code repair with unit tests,
formal theorem proving, exact-answer reasoning, database queries, or games with
a checkable terminal state, have an objective stopping criterion and often admit
structured solution paths. In this setting, using one known successful trace as
a reference is a useful approximation for accounting: it is not assumed to be
the only possible way to solve \(x\), but it gives a concrete baseline for
measuring tool selection, detours, and resource waste. Successes obtained
through other valid traces are possible and are treated separately below.

In the reference-trace regime, a generation trial verifies only when it produces
\(\tau^\star(x)\). Conditional on producing this reference trace, verification
succeeds with probability \(p(x)>0\). Let \(q^\star(x,z)\) be the probability,
under the tool-allocation policy \(q(\cdot\mid z,\tau_{<k})\), of producing
\(\tau^\star(x)\). Define
\[
t_0(x):=\frac{1}{p(x)},
\]
the expected number of trials needed to verify if the reference trace were
produced on every trial. Empirically, \(p(x)\) can be estimated by repeatedly
executing the agent while forcing, replaying, or conditioning on
\(\tau^\star(x)\), and measuring the fraction of trials that verify.

\begin{proposition}[Tool-allocation first-passage identity]
\label{prop:tool-first-passage}
In the reference-trace regime, the per-trial verification probability is
\[
\Prob(V(x,Y)=1\mid x,z)
=
q^\star(x,z)p(x).
\]
If each trial has conditional mean cost \(\overline k_\sigma(x,z)\), and any signal cost
is zero or already included in the trial cost, then
\begin{equation}
\E\!\left[B^{\mathrm{hit}}(x,z)\mid x,z\right]
=
\overline k_\sigma(x,z)\,
\frac{t_0(x)}{q^\star(x,z)},
\qquad
t_0(x):=\frac{1}{p(x)}.
\label{eq:tool-first-passage}
\end{equation}
\end{proposition}

We now decompose the tool-allocation penalty. Across deployment instances,
reference traces may share latent solution modes: groups of instances that can
be solved by a common solution pattern, such as a similar reasoning structure,
tool-use pattern, or sequence of external calls. For each instance \(x\), let
\(\tau^\star(x)=(C_1^\star,\ldots,C_{L(x)}^\star)\) be the selected reference
tool-use trace. At step \(k\), define
\begin{equation}
\rho_k(c\mid z,\tau_{<k})
:=
\Prob\!\left(
C_k^\star=c
\,\middle|\,
Z=z,\tau^\star_{<k}=\tau_{<k},L\ge k
\right).
\label{eq:reference-batch-posterior}
\end{equation}
This is an oracle distribution over the next reference batch of tool calls. It
is one-hot when the next batch is fixed by the reference trace, and
non-degenerate when several successful reference continuations remain possible
under the same signal and prefix. Assume
\(\rho_k(c\mid z,\tau_{<k})>0\) implies
\(q_\sigma(c\mid z,\tau_{<k})>0\). Define
\begin{align}
G_Z
&:=\sum_{k\ge1}\Prob(L\ge k)
I(C_k^\star;Z\mid \tau_{<k}^\star,L\ge k),\nonumber\\
\varepsilon_\sigma
&:=\sum_{k\ge1}\Prob(L\ge k)
\E\!\left[
\KL\!\left(
\rho_k(\cdot\mid Z,\tau_{<k}^\star)
\Vert
q_\sigma(\cdot\mid Z,\tau_{<k}^\star)
\right)
\,\middle|\,L\ge k\right].
\label{eq:batched-tool-mismatch}
\end{align}

The term \(G_Z\) is the useful information supplied by the signal about which
next tool batch should be used. The only mismatch term is \(\varepsilon_\sigma\):
it is small when the agent assigns high probability to the reference next batch,
and large when it calls the wrong tools, calls them in the wrong order, or
allocates probability to irrelevant batches. Operationally, \(\rho_k\) is
estimated from benchmark or teacher reference traces, while \(q_\sigma\) is
estimated by repeated generation trials of the same agent system on the same
instance or on instances from the same latent solution mode.

\begin{lemma}[Batched tool-allocation gap]
\label{lem:tool-cross-entropy}
For a system \(\sigma\) using signal \(Z\), write
\[
\Xi(\sigma,Z)
:=
\E\!\left[-\sum_{k=1}^{L}
\log q_\sigma(C_k^\star\mid Z,\tau_{<k}^\star)\right].
\]
If two systems \(\sigma_0,\sigma\) share the same deployment distribution and
reference call batches, then
\begin{equation}
\Xi(\sigma_0,Z_0)-\Xi(\sigma,Z)
=G_Z-G_{Z_0}+\varepsilon_{\sigma_0}-\varepsilon_\sigma.
\label{eq:tool-cross-entropy}
\end{equation}
Thus \(G_Z\) measures how much the signal helps identify appropriate tool
batches, while \(\varepsilon_\sigma\) is the single tool-policy mismatch term.
The intrinsic diversity of reference tool-use patterns is common to both
systems and cancels in this relative accounting.
\end{lemma}

Define the expected-log conditional first-passage diagnostic
\begin{equation}
\mathcal L(\sigma):=\E_{X,Z}\!\left[
\log\E(B_\sigma^{\mathrm{hit}}\mid X,Z)\right].
\label{eq:tool-risk}
\end{equation}
It is neither \(\E\log B_\sigma^{\mathrm{hit}}\),
\(\log\E B_\sigma^{\mathrm{hit}}\), nor \(\log T_\sigma(\delta)\).

\begin{theorem}[Four-term tool-allocation decomposition]
\label{thm:four-term}
Let agent systems \(\sigma_0,\sigma\) share the task distribution, tool set,
and reference tool combinations, and use signals \(Z_0,Z\).  Suppose both satisfy the
reference-combination regime and have no separate signal cost.  Define
\begin{align*}
C&:=\E_{X,Z_0}\log\overline k_{\sigma_0}(X,Z_0)
-\E_{X,Z}\log\overline k_\sigma(X,Z),\\
\Phi&:=\E_X\!\left[\log\frac{t_0(\sigma_0,X)}{t_0(\sigma,X)}\right].
\end{align*}
Then
\begin{equation}
\mathcal L(\sigma_0)-\mathcal L(\sigma)
=C+\Phi+G_Z-G_{Z_0}+\varepsilon_{\sigma_0}-\varepsilon_\sigma.
\label{eq:four-term-general}
\end{equation}
If \(Z_0\) is conditionally uninformative about every reference dispatch
and the baseline policy matches the corresponding prior, then
\(G_{Z_0}=\varepsilon_{\sigma_0}=0\), and
\begin{equation}
\boxed{\Delta_{\mathrm{tool}}=C+\Phi+G_Z-\varepsilon_\sigma.}
\label{eq:four-term-baseline}
\end{equation}
The four terms are relative mean trial cost, focused competence, signal
information about appropriate tool use, and batched tool-policy mismatch.
Because the policy can change the trial-cost distribution, this is a
mechanistic accounting rather than a causal separation.
\end{theorem}

\begin{proposition}[Testable off-reference departure]
\label{prop:off-reference}
Let
\[
r_\sigma^{\mathrm{alt}}(x,z)
:=\Prob(V(x,Y_r)=1,\mathbf C\ne\mathbf C^\star(x)\mid x,z)
\]
be the per-retry verified-success probability contributed by non-reference
executions.  Ignoring signal cost, define the log residual as the actual
conditional mean minus the reference-combination prediction:
\begin{equation}
r_\sigma(x,z)=-\log\!\left(
1+\frac{r_\sigma^{\mathrm{alt}}(x,z)}{q_\sigma^\star(x,z)p_\sigma(x)}
\right).
\label{eq:off-reference-residual}
\end{equation}
If the ratio inside the logarithm is at most \(\xi\), then
\(|r_\sigma(x,z)|\le\log(1+\xi)\).  A separate signal cost breaks the exact
log-additivity.  The verifier still accepts every valid non-reference
execution; the proposition only tests the adequacy of the diagnostic regime.
\end{proposition}

Separating \(\overline k_\sigma\) from \(t_0\) requires trial-level resource and
verification logs.  A policy change can alter both mean trial cost and
tool-policy mismatch, so the terms are not causal without randomized
interventions.  Proper time remains the model-selection criterion and is
measured directly; Theorem~\ref{thm:four-term} audits why expected log
conditional first-passage resource differs when its assumptions are adequate.

\begin{algorithm}[How to Pick the Model]
\label{alg:pick-model}
\begin{tabular}{@{}r p{0.86\linewidth}@{}}
\textit{Input:} &
Verifiable task, candidate agent systems \(\Sigma\), verifier \(V\),
resource unit, tolerance \(\delta\), horizon \(B\),
and confidence level \(1-\alpha\).\\[0.25em]
\textit{Output:} &
A statistically certified winner or ``no winner,'' plus an optional batched
tool-allocation mechanism audit.\\[0.35em]
1 & \textit{Freeze the target.} Declare \(V\), the resource unit, \(\delta\),
\(B\), and the candidate menu.\\
2 & \textit{Measure first passage.} Stop each run at its first verified artifact
or at \(B\); retain failures as infinite and pair system runs by task instance.\\
3 & \textit{Construct bands.} Compute Equation~\eqref{eq:empirical-cdf}, the
simultaneous DKW bands, and \([T_\sigma^-,T_\sigma^+]\) for every system.\\
4 & \textit{Certify or abstain.} Name a winner only if it is admissible and its
upper proper-time bound lies below every competitor's lower bound.\\
5 & \textit{Audit the mechanism.} On disjoint calibration and confirmation data,
estimate reference call-batch probabilities, conditional success, variable
retry costs, and any incremental signal cost; test IID retries,
reference-combination adequacy, and off-reference success.\\
6 & \textit{Explain, do not rerank.} Where the reference-combination regime is
adequate, report the four terms, batched tool-policy mismatch, and held-out
residual; keep the direct proper-time certificate as selection authority.
\end{tabular}
\end{algorithm}

The chain is explicit: Definition~\ref{def:proper-time} sets the deployment
estimand; Corollary~\ref{thm:finite-sample} selects or abstains;
the retry ledger links variable trial costs to mean first passage; and
Theorem~\ref{thm:four-term} supplies a scoped tool-policy
mechanism audit.
\label{sec:theory-end}
